\section{Purpose and Standard of Success}

\textbf{Why this matters:} Without a stated standard, "does this work" has no answer, and a system ends up judged against whatever comparator is convenient --- usually an overworked person, which makes almost anything look like an improvement.

\textbf{Practical implication:} Before evaluating a tool, agree explicitly on what a strong human comparator actually achieves, and require tail risks to be controlled separately from the average. A good average with an unacceptable tail is not a pass.

Everything from here forward is a catalogue of what can go wrong, why verification is hard, how errors propagate, and what an organization needs in place to catch and contain them. Read on its own, that catalogue can look like an argument for never deploying anything --- a wall of failure modes with no stated finish line. That reading would be a mistake, and it's worth stating plainly what this framework is actually for.

\textbf{The purpose is not to demand an error-free system.} No human process and no technical system offers perfect availability, complete information, or uniformly correct performance under every possible condition --- that standard has never been met by anything, including the all-human processes AI is being compared against, and holding AI to it uniquely would block useful deployment without producing any real gain in safety. The right question is never \emph{"can this workflow produce an error"} --- every useful system can --- but rather: \textbf{under a defined range of operating conditions, does the complete human-AI workflow produce fewer and less consequential errors than a well-rested, properly trained human performing the same work, while also improving speed, capacity, or cost?}

\textbf{The comparison has to be made against a strong human baseline, not a convenient one.} Measuring an AI system against an inexperienced, fatigued, or overloaded person sets an artificially low bar and will make almost anything look like an improvement. Reduced fatigue and higher throughput are real, legitimate benefits of automation --- but they should count as benefits \emph{on top of} competent baseline performance, not as a substitute for actually reaching it. A system that merely outperforms an overworked, undertrained comparator hasn't demonstrated much; establishing what a genuinely well-resourced human process actually achieves is real, sometimes uncomfortable work in its own right, and a comparison is only as good as the honesty of that baseline.

\textbf{Availability needs to be split into at least four distinct questions, because collapsing them makes a system look more reliable than it is.} A system can be technically available --- it responds, promptly, without error --- while being \emph{epistemically} unavailable: it lacks the documents, the current law, the permissions, or the actual competence the task requires, and produces a fluent answer anyway. Worth distinguishing explicitly: \textbf{service availability} (does the system run at all); \textbf{task availability} (can it competently perform \emph{this specific} task, as opposed to tasks in general); \textbf{assurance availability} (is there actually enough evidence to justify relying on this particular result, per the verification-strength distinctions in Section 7 and the assurance-argument treatment in Section 13); and \textbf{fallback availability} (can the work continue safely through some other route when the system can't do it). This matters for a specific reason: \textbf{an abstention should not automatically be counted as an error.} A system that declines, asks a clarifying question, or escalates to a human when it lacks sufficient grounding --- the well-behaved response to the "environment" and "groundedness" failures described in Section 9 --- is often a successful control acting exactly as it should, not a failure to complete the task. The caveat runs the other way too: a system that abstains too readily becomes useless regardless of how safe it is, so abstention is a dial to calibrate deliberately against the task's actual stakes, not a free action to maximize.

\textbf{Harm has to be measured directly, not approximated by an error count.} A workflow can have fewer total errors and still be unacceptable, if the errors that remain are unusually severe, correlated, or hard to detect --- exactly the consequence dimensions named in Section 13 (severity, detectability, reversibility, blast radius, persistence). The right comparison needs two things held simultaneously, not one averaged into the other: \textbf{common-case performance} (fewer errors, faster completion, lower reviewer effort across the bulk of ordinary, representative work) and \textbf{tail-risk constraints} (specific categories of outcome that must stay below a defined threshold \emph{regardless} of how good the average looks). A contract-review system might substantially reduce ordinary issue-spotting mistakes while still being unfit for deployment if it occasionally crosses a matter barrier or silently drops an entire document --- a good average with an intolerable tail is not a passing result, and expected-loss arithmetic that blends the two into one number will hide exactly the failure mode that matters most.

\textbf{"Most circumstances" is directionally right but needs to become a defined, tested operating envelope, or it becomes a way to quietly exclude the hard cases until the system looks better than it is.} A real envelope specifies: the supported tasks and document types; the jurisdictions, languages, and subject areas actually covered; the input quality assumed; the data and tool dependencies required; the user competence the deployment assumes; the maximum workflow complexity handled; the circumstances that require escalation rather than autonomous completion; and --- just as importantly --- the known unsupported or weakly supported cases, stated explicitly rather than discovered by accident. Within that envelope, the ambition can legitimately be a workflow materially more reliable than a strong human comparator. Outside it, the standard shifts to \textbf{safe degradation}: decline, request the missing information, narrow the conclusion to what's actually supported, escalate, or hand the task back to a human-controlled process entirely. A system is not unreliable merely because its competence has real boundaries --- every system's does. It becomes unreliable specifically when it cannot recognize where those boundaries sit, or cannot enforce them once it's outside them.

\textbf{The comparison that actually matters is workflow versus workflow, not model versus human.} Benchmarking a raw model against an unassisted person understates what's actually being deployed and overstates what a bare model result tells you. The comparison worth running spans several configurations: a trained human using conventional tools alone; a trained human working with the AI system; the AI system operating under mandatory human review; the AI system operating with escalation rules and deterministic checks in place; and, only where it's genuinely been earned, the AI system operating autonomously within a narrow, tested envelope. The frequent and often underappreciated result is that \textbf{human-plus-system outperforms either the human or the system alone} --- which is itself a reason the comparison has to include that combined configuration explicitly rather than treating "human" and "AI" as the only two options. The metrics worth tracking go well past raw accuracy: consequential error rate specifically (not just any error), omission rate, detection latency, correction and rework required, throughput, time per completed task, reviewer effort, abstention and escalation rates, consistency across equivalent cases, severe near misses, performance under degraded conditions, and --- a subtler one --- performance over time as users grow accustomed to the system, since habituation itself changes the workflow's actual reliability in ways a one-time evaluation won't show.

\textbf{Error budgets are a natural extension of "availability" thinking, but only if they're allowed to be multidimensional.} A single permissible error percentage is close to meaningless here, since it treats a harmless formatting slip and an unauthorized disclosure as the same unit. A workable budget instead allocates tolerance by failure class and consequence: a modest allowance for harmless formatting defects; a smaller one for substantive analytical errors; a very small one for anything client-facing; and, for a defined set of failure classes --- an unauthorized disclosure, an uncontrolled irreversible action --- effectively zero tolerance regardless of how good the aggregate success rate looks. Some errors are genuinely economically tolerable as a cost of doing useful work faster. Others warrant immediate suspension of the workflow no matter what the overall numbers say, and deciding which is which in advance, by failure class, is what makes an error budget a real operational tool rather than a single misleading average.

Put together, this gives the framework a positive standard rather than only a list of hazards to guard against:

\begin{quote}
The purpose of this framework is not to demand an error-free AI system. No human or technical system can provide perfect availability, complete information, or uniformly correct performance under every possible condition. A requirement of perfection would prevent useful deployment without meaningfully improving safety.

The appropriate standard is comparative and operational. An AI-enabled workflow should be adopted where, within a clearly defined and tested operating envelope, the complete system --- including its human users, data, models, controls, interfaces, verification mechanisms, and fallback procedures --- performs the relevant work more efficiently and with materially fewer or less consequential errors than a well-rested, properly trained human using the best reasonably available conventional process.

That comparison must not rest on average accuracy alone. The system must also control severe and difficult-to-detect failures, preserve required confidentiality and authority boundaries, recognize when it is outside its competence, and degrade safely when necessary information, infrastructure, or assurance is unavailable. A system that improves ordinary performance while introducing an unacceptable class of rare, irreversible, or correlated failures has not met the standard.

The objective is therefore not maximum automation. It is the greatest justified improvement in reliability and efficiency, subject to explicit constraints on consequence, detectability, reversibility, and residual risk. Human involvement should remain wherever it produces a material improvement in those outcomes, and should be reduced only where evidence shows automation performs the function more reliably.

\end{quote}

Which gives the whole framework a single closing test to apply to any specific deployment, rather than an open-ended list of concerns to weigh informally: \emph{does this system, within its established operating envelope, outperform competent human practice on efficiency and total risk --- and does it fail safely outside that envelope?} That's demanding, in that it requires actually establishing the comparator, the envelope, and the tail-risk constraints rather than asserting them --- but it is not a demand for perfection, and every section that follows should be read as filling in what's needed to answer that one question honestly, not as an independent reason to hesitate.
